% Generated draft: requires mathematical and source review.
\documentclass[11pt]{article}
\usepackage[margin=25mm]{geometry}
\usepackage{fontspec}
\setmainfont{Latin Modern Roman}
\usepackage{amsmath,amssymb,mathtools,unicode-math}
\setmathfont{Latin Modern Math}
\usepackage{enumitem,needspace,xurl,hyperref}
\setlength{\parindent}{0pt}
\setlength{\parskip}{5pt}
\setlength{\emergencystretch}{3em}
\newcommand{\N}{\mathbb{N}}
\newcommand{\Z}{\mathbb{Z}}
\newcommand{\Q}{\mathbb{Q}}
\newcommand{\R}{\mathbb{R}}
\newcommand{\C}{\mathbb{C}}
\newcommand{\F}{\mathbb{F}}
\newcommand{\A}{\mathbb{A}}
\newcommand{\PP}{\mathbb{P}}
\newcommand{\E}{\mathbb{E}}
\newcommand{\eps}{\varepsilon}
\newcommand{\dd}{\,\mathrm{d}}
\begin{document}
\section*{UnsolvedMath ID: 30000334}
\textbf{Source problem code:} catalogue entry
\subsection{Definitions and mathematical statement}
The source asks for an effective proof of an existing theorem of Thurston, not for a new assertion of qualitative fixed-point existence. In the usual setting, a suitable hyperbolizable three-manifold \(M\) with marked boundary has a skinning map
\[
\sigma_M:\mathcal T(\partial M)\longrightarrow\mathcal T(\overline{\partial M}),
\]
where \(\mathcal T\) denotes Teichm\"uller space and the bar reverses orientation. The map records the other conformal boundary data of appropriate boundary-subgroup covers of the geometrically finite structure. Exact manifold, boundary, marking, and deformation-space hypotheses are not reproduced in the record.

To state a fixed-point problem, one generally also uses the source's gluing or orientation-reversing identification, inducing a map \(\iota:\mathcal T(\overline{\partial M})\to\mathcal T(\partial M)\). The resulting self-map is \(F=\iota\circ\sigma_M\), with the appropriate product convention if there are several pieces. The qualitative target has the form
\[
\exists X\in\mathcal T(\partial M)\qquad F(X)=X
\]
under the hypotheses of the intended Thurston theorem. A map into the oppositely oriented Teichm\"uller space is not silently treated as a self-map without these identifications.

The requested strengthening is to make the proof effective. This could mean computable a priori bounds from the manifold and gluing data, a certified approximation procedure for a fixed point to arbitrary accuracy, or quantitative control of the proof's limiting or contraction arguments. The source does not define which notion of effectivity, input encoding, output certification, or error metric is required. It also does not supply the hypotheses excluding gluing obstructions. Therefore no particular algorithmic theorem, polynomial-time requirement, or unrestricted fixed-point assertion for all three-manifolds is invented.

The supplied assessment checked 2026-08-21 reports that later effective bounds do not remove the general request for an effective proof. It does not reproduce those bounds or the exact remaining effectiveness gap.
\subsection{Short English statement}
Make Thurston's skinning-map fixed-point proof effective under its intended geometric hypotheses. The relevant self-map includes gluing identifications, and the source does not specify the quantitative or algorithmic meaning of effective.
\subsection{Sources}
\begin{itemize}
\item Local UnsolvedMath registry, permanent ID 30000334, catalogue entry, \emph{Effective Thurston Fixed-Point Theorems for Skinning Maps}.
\item \emph{Low-Dimensional Manifolds} (2006), Oberwolfach Reports, pp. 2519--2570, DOI \url{https://doi.org/10.4171/owr/2005/45}; the fixed-point theorem is attributed to Thurston in the supplied question.
\item Supplied reference labeled ``Skinning maps survey'': \url{https://people.math.wisc.edu/~aekent2/bartholomew.pdf}.
\item Supplied literature triage checked 2026-08-21; no exact theorem statement, effective bound, or full survey citation is reproduced.
\end{itemize}
\end{document}
